\naslov{Linearne NDE višjega reda}

Obravnavamo enačbe oblike
\[
  a_n(t) x^{(n)} + \ldots + a_1(t) \dot{x} + a_0(t) x = b(t).
\]

\begin{definicija}
  Linearni diferencialni operator s koeficienti $a_i(t)$ je preslikava
  \[
	L : \zvodv{1}{[a,b]} \to \zvezne{[a,b]},
  \]
  podana s predpisom
  \[
	Lx(t) = a_n(t) x^{(n)} + \ldots + a_0(t) x.
  \]
\end{definicija}

Splošna rešitev homogene enačbe $Lx = 0$ je $\jedro L$.
Vemo, da je splošna rešitev $n$-dimenzionalni vektorski prostor.
Enačbo s substitucijo
\begin{align*}
  x_1 &= x \\
  x_2 &= \dot{x} \\
  \vdots & \\
  x_n &= x^{(n-1)}
\end{align*}
prepišemo v sistem
\begin{gather*}
  \dot{x}_0 = x_1 \\
  \dot{x}_1 = x_2 \\
  \vdots \\
  \dot{x}_n = \inv{a_n} (b - a_0 x_1 - a_2 x_1 - \ldots - a_{n-1} x_{n-2})
\end{gather*}
Označimo
\[
  p_i(t) = \frac{a_i(t)}{a_n(t)},
\]
s čimer izrazimo matriko koeficientov zgornjega sistema
\[
  A(t) =
  \begin{bmatrix}
	0 & 1 & 0 & \cdots & 0 \\
	0 & 0 & 1 & \cdots & 0 \\
	0 & 0 & 0 & \cdots & 0 \\
	\vdots & \vdots & \vdots & \ddots & \vdots \\
	-p_0 & -p_1 & -p_2 & \cdots & -p_{n-1}
  \end{bmatrix}.
\]

\begin{izrek}
  Naj bo $L: \zvodv{n}{[a,b]} \to \zvodv{0}{[a,b]}$ regularen diferencialni
  operator.
  Za vsak začetni pogoj $x(t_0) = c_0, \dot{x}(t_0) = c_1, \ldots,
  x^{(n-1)}(t_0) = c_{n-1}$ ima enačba $Lx = b$ natanko eno rešitev na vsem
  intervalu $[a,b]$.
\end{izrek}

\begin{opomba}
  Diferencialni operator: $Lx = a_n x^{(n)} + \ldots + a_1 \dot{x} +
  a_0 x = b$ je regularen, če je $a_n(t) \ne 0$ za vsak $t$ in če so $a_i(t)$
  omejene.
\end{opomba}

Množica rešitev homogene linearne enačbe $Lx = 0$ je $n$-dimenzionalen vektorski
prostor v $\zvodv{n}{[a,b]}$.
Vsaka rešitev $x(t)$ namreč na enoličen način določa rešitev sistema
$\dot{\vec{x}} = A \vec{x}$ za
\[
  \dot{\vec{x}}
  =
  \begin{bmatrix}
	x(t) \\ \dot{x}(t) \\ \vdots \\ x^{(n-1)}(t)
  \end{bmatrix}.
\]
Tudi obratno je res: vsak vektor $\vec{x}$ na enoličen način določa svojo prvo
komponento.

Oglejmo si fundamentalno matriko
\[
  \phi(t)
  =
  \begin{bmatrix}
	x_1 & x_2 & \cdots & x_n \\
	\dot{x}_1 & \dot{x}_2 & \cdots & \dot{x}_n \\
	\vdots & \vdots & \ddots & \vdots \\
	x_1^{(n-1)} & x_2^{(n-1)} & \cdots & x_n^{(n-1)}
  \end{bmatrix}.
\]
Denimo, da so funkcije $x_1(t), x_2(t), \ldots, x_n(t)$ baza rešitev enačbe $Lx
= 0$.
Za bazo velikokrat vzamemo take vektorje $\vec{x}_1, \ldots, \vec{x}_n$, da je
$\phi(t=0) = I$.
Če je ta baza dobljena iz baze rešitev enačbe $Lx = 0$, potem za te rešitve
velja $x_i(0) = 0, \ldots, x_i^{(i-1)}(t) = 0, x_i^{(i)}(0) = 1, x_i^{(i+1)}(0)
= 0, \ldots, x_i^{(n-1)}(0) = 0$.
Determinanta
\[
  W(t) = \det \phi(t)
\]
se imenuje \pojem{determinanta Wronskega}.
V tem primeru se Liouvilleova formula glasi
\[
  W(t) = W(t_0) \exp\left( - \int_{t_0}^t p_{n-1}(\tau) d\tau \right).
\]

Rešitev nehomogene enačbe dobimo s pomočjo variacije konstante;
\[
  \vec{x} = \int_{t_0}^t \phi(t) \phi^{-1}(\tau) \vec{b}(\tau) d\tau,
\]
oziroma, ker ima $\vec{b}$ v tem primeru le eno neničelno komponento $b(t)$, bo
prva komponenta $\vec{x}$ enaka
\[
  x(t)
  = \int_{t_0}^t \sum_{i=1}^n \phi_{1i}(t) \phi_{in}^{-1}(\tau) b(\tau) d\tau
  = \sum_{i=1}^n x_i(t) \int_{t_0}^t \phi_{in}^{-1}(\tau) b(\tau) d\tau,
\]
kjer je $(x_i(t))_i$ baza rešitev homogene enačbe $Lx = 0$.

\vprasanje{Kako izračunaš rešitev linearne NDE višjega reda?}

\podnaslov{Enačbe s konstantnimi koeficienti}

Naj bo sedaj linearen diferencialni operator $L$ podan z
\[
  Lx = x^{(n)} + a_1 x^{(n-1)} + \ldots + a_n x.
\]
Oglejmo si enačbo $Lx = 0$.
Če vstavimo nastavek $x(t) = e^{\lambda t}$:
\[
  \lambda^n e^{\lambda t} + a_1 \lambda^{n-1} e^{\lambda t} + \ldots + a_{n-1}
  \lambda e^{\lambda t} + a_n e^{\lambda t} = 0
\]
oziroma (ker $e^{\lambda t} \ne 0$ tudi za $\lambda \in \C$)
\[
  \lambda^n + a_1 \lambda^{n-1} + \ldots + a_{n-1} \lambda + a_n = 0.
\]
Ta polinom imenujemo \pojem{karakteristični polinom enačbe $Lx = 0$} in označimo
s $P(\lambda)$.

\begin{trditev}
  Če so $\lambda_1, \ldots, \lambda_n$ različne ničle karakterističnega polinoma
  $P(\lambda)$, potem so funkcije $x_i(t) = e^{\lambda_i t}$ baza rešitev
  homogene enačbe $Lx = 0$.
\end{trditev}

\begin{proof}
  Vemo, da je rešitev sistema vektorski prostor, dokazati moramo samo, da so te
  rešitve linearno neodvisne.
  Priredimo našim rešitvam pripadajoče rešitve sistema, ki je prirejen $Lx = 0$,
  \[
	x_i(t) \mapsto \vec{x}_i(t) =
	\begin{bmatrix}
	  x_i(t) \\ \dot{x}_i(t) \\ \vdots \\ x_i^{n-1}(t)
	\end{bmatrix}.
  \]
  Te stolpce zložimo v matriko in dobimo kandidatko za fundamentalno matriko
  $\phi(t)$
  \[
	\phi(t)
	=
	\begin{bmatrix}
	  e^{\lambda_1 t} & e^{\lambda_2 t} & \cdots & e^{\lambda_n t} \\
	  \lambda_1 e^{\lambda_1 t} & \lambda_2 e^{\lambda_2 t} & \cdots & \lambda_n
																	   e^{\lambda_n t} \\
	  \lambda_1^2 e^{\lambda_1 t} & \lambda_2^2 e^{\lambda_2 t} & \cdots & \lambda_n^2
																	   e^{\lambda_n t} \\
	  \vdots & \vdots & \ddots & \vdots \\
	  \lambda_1^{n-1} e^{\lambda_1 t} & \lambda_2^{n-1} e^{\lambda_2 t} & \cdots
												 & \lambda_n^{n-1} e^{\lambda_n t} \\
	\end{bmatrix}
  \]
  Matrika $\phi(0)$ je vandermondova, torej
  \[
	W(t) = \det \phi(t) = W(0) \exp\left( \int_0^t \operatorname{sl}(A) d\tau
	\right)
	= \prod_{i>j} (\lambda_i - \lambda_j) \cdot e^{- a_1 t}.
  \]
  Ker so vsi $\lambda_i$ različni, velja $W(0) \ne 0$, torej je $W(t) \ne 0$
  za vsak $t$, in so funkcije $x_i$ res linearno neodvisne in so baza prostora
  rešitev enačbe $Lx = 0$.
\end{proof}

\vprasanje{Kaj je karakteristični polinom homogene linearne NDE višjega reda?
  Kako z njim poiščemo rešitve enačbe? Dokaži.}

\podnaslov{O karakterističnemu polinomu}

Z razvojem po prvem stolpcu lahko izračunamo
\[
  \det (A - \lambda I)
  =
  \begin{vmatrix}
	-\lambda & 1 & 0 & \cdots & 0 \\
	0 & -\lambda & 1 & \cdots & 0 \\
	0 & 0 & -\lambda & \cdots & 0 \\
	\vdots & \vdots & \vdots & \ddots & \vdots \\
	-a_n & -a_{n-1} & -a_{n-2} & \cdots & -a_1-\lambda
  \end{vmatrix}
  = \lambda^n + a_1 \lambda^{n-1} + \ldots + a_{n-1} \lambda + a_n.
\]
Lastne vrednosti matrike $A$ so torej res ničle karakterističnega polinoma
$P(\lambda)$ enačbe $Lx = 0$.

% LocalWords:  Wronskega Liouvilleova vandermondova
